\documentclass[10pt,a4paper]{amsart}
\usepackage[margin=19mm]{geometry}
\usepackage{amsmath,amssymb,mathtools}
\usepackage[scr=rsfs,cal=euler]{mathalfa}
\usepackage{xfrac}
\usepackage[unicode,hidelinks,bookmarks=false]{hyperref}
\newcommand{\ie}{i.e.\ }
\newcommand{\cf}{cf.\ }
\newcommand{\resp}{respectively\ }
\newcommand{\Subsection}{\subsection}
\providecommand{\nobreakdash}{\nobreak\hbox{-}}
\setlength{\emergencystretch}{2em}
\multlinegap0pt
\usepackage{bm}
\usepackage{bbm}
\usepackage{dsfont}
\usepackage{xspace}
\usepackage{cancel}
\usepackage{mathtools}
\usepackage{amsthm}
\newtheorem{proposition}{Proposition}
\title[An approximate counterexample to the Barker--Larman problem ]{An approximate counterexample to the Barker--Larman problem in dimension $4$}
\author{J. Haddad}
\date{Source: arXiv:2609.10257v2 (CC BY 4.0)}
\begin{document}
\maketitle

\noindent\emph{Source and licence.} An approximate counterexample to the Barker--Larman problem in dimension $4$. Version arXiv:2609.10257v2 (CC BY 4.0), distributed under the Creative Commons Attribution 4.0 International licence, as recorded in the arXiv licence metadata for this exact version and shown on the abstract page https://arxiv.org/abs/2609.10257v2. Full rights notice: licenses/arxiv-2609.10257-CC-BY-4.0.txt.\par\smallskip

\noindent\textit{Editorial scope.} Counted unit: the Proposition whose source label is \texttt{prop\_bound\_Rr} in the article (source0.tex lines 730--735, source label \texttt{prop\_bound\_Rr}), delivered together with the article's own complete original proof of it (source0.tex lines 736--762, one \texttt{proof} environment of 27 lines). No mathematics is written, repaired or re-derived: statement and proof are the article's own text line for line. Required context retained uncounted: eq\_parametrization\_gamma (lines 187--191); eq\_varphi1 (lines 476--479); eq\_series\_bound (lines 679--682). Every definition, display and label of those lines is retained unchanged. Omitted from this package and not redistributed: 1--186, 192--475, 480--678, 683--729, 763--898. Nothing inside a retained excerpt is omitted or altered: every retained body is delivered exactly as the article's lines are written, so no omission receipt of any kind accompanies this package. Citations: the retained excerpts carry no citation command at all, so no bibliography entry of the article is transcribed and no reference is redistributed. Reference adaptation: none was needed and none was made; every reference inside the delivered unit resolves inside it, so no unresolved reference or citation is produced. Printed slips kept literal: no printed word, symbol, hypothesis or number of the counted statement or of its proof was altered. Layout and class: the article's own class is \texttt{amsart} with packages including hyperref, verbatim, nccmath, graphicx, enumitem, mathtools; the wrapper is the lane's pinned \texttt{amsart} editorial wrapper at 10pt on A4, which restates the theorem-like environments the retained text needs and adds the article's own macro definitions that the retained text needs (none), with extra wrapper packages (bm, bbm, dsfont, xspace, cancel, mathtools). The delivered numbering is the unit's own. Assets: no retained span contains a figure environment, a graphics-inclusion command, a PDF-page inclusion, a table or a TikZ picture; no image file is copied into this package, the package declares no asset dependency and no image file is redistributed with it. Licence: version arXiv:2609.10257v2 of this article is distributed under the Creative Commons Attribution 4.0 International licence, as recorded in the arXiv licence metadata for this exact version and shown on the abstract page https://arxiv.org/abs/2609.10257v2. A full rights notice is delivered at licenses/arxiv-2609.10257-CC-BY-4.0.txt. Characters: every retained excerpt is pure ASCII, so the delivered engine, XeLaTeX with its default Latin Modern text font, reports no missing character.
	\begin{align}
        \label{eq_parametrization_h_plus} \delta_{K,+}(x) &= (\cos(x), \sin(x)) +\tan\left(h^{-1}(x) - x\right) (\cos(x+\pi/2), \sin(x+\pi/2)),\\
	\label{eq_parametrization_h_minus} \delta_{K,-}(x) &= (\cos(x), \sin(x)) +\tan\left(h^{-1}(-x) + x\right) (\cos(x-\pi/2), \sin(x-\pi/2)),\\
	\label{eq_parametrization_gamma} \gamma_K(t) &= \cos(t-h(t))^{-1} (\cos(t), \sin(t)),
	\end{align}

\begin{equation}
	\label{eq_varphi1}
	\varphi_1(t) = -\frac 1{48 \pi} + \frac 1{24 \pi} (1+2\cos(2t)).
\end{equation}

\begin{equation}
	\label{eq_series_bound}
	\sum_{n=1}^N \frac{r^n}{n!} \|\varphi_n'\|_\infty \leq \frac 12 \text{ and } \sum_{n=1}^N \frac{r^n}{n!} \|\varphi_n\|_\infty \leq \frac \pi {12}.
\end{equation}

\begin{proposition}
	\label{prop_bound_Rr}
	For every $N>0$ there is $r_N$ such that for $\lambda \in (\lambda_0-r_N, \lambda_0+r_N)$,
	\[ |R(K_{\lambda,N}) - r(K_{\lambda,N})| \geq c |\lambda - \lambda_0|, \]
	where $c>0$ is a universal constant.
\end{proposition}

\begin{proof}
	Let $r_N>0$ be small enough so that \eqref{eq_big_angle} and \eqref{eq_small_angle} hold for all $\lambda \in (\lambda_0-r_N,\lambda_0+ r_N)$, and such that \eqref{eq_series_bound} holds with $r = r_N$.
	According to our parametrization \eqref{eq_parametrization_gamma} of $K_{\lambda,N}$, for each $t \in [0,\pi]$ and $\lambda \in (\lambda_0-r_N, \lambda_0 + r_N)$, there is a point $\gamma_{K_{\lambda,N}}(t)$ at distance $\cos(h(t)-t)^{-1}$ from the origin.
	Consider only $\lambda \in [\lambda_0, \lambda_0 + r_N)$. Taking a smaller $r_N$ if necessary, we may bound $h_{\lambda,N}$ using only the linear part,
	\begin{align}
		\label{eq_big_angle}
		-h_{\lambda,N}(0) &= \frac \pi 3 - (\lambda-\lambda_0) \varphi_1(0) - \sum_{n=2}^N \frac {(\lambda-\lambda_0)^n}{n!} \varphi_n(0)\\
		&\leq \frac \pi 3 - \frac 12 (\lambda-\lambda_0) \varphi_1(0)\\
		&= \frac \pi 3 - \frac 5{96 \pi} (\lambda-\lambda_0),
	\end{align}
	and
	\begin{align}
		\label{eq_small_angle}
	\frac \pi 2 - h_{\lambda,N}\left(\frac \pi 2\right)
		&= \frac \pi 3 - (\lambda-\lambda_0) \varphi_1(\pi/2) - \sum_{n=2}^N \frac {(\lambda-\lambda_0)^n}{n!} \varphi_n(\pi/2)\\
		&\geq \frac \pi 3 - \frac 12 (\lambda-\lambda_0) \varphi_1(\pi/2)\\
		&= \frac \pi 3 + (\lambda-\lambda_0) \frac 1{32 \pi},
	\end{align}
	where we used \eqref{eq_varphi1} in both inequalities, and the fact that $\varphi_1(\pi/2) < 0 < \varphi_1(0)$.

	Taking $t = 0, \frac \pi 2$ in \eqref{eq_parametrization_gamma} we obtain for $\lambda \in [\lambda_0, \lambda_0 + r_N),$
	\[R(K_{\lambda,N}) \geq \cos\left(\frac \pi 2 - h_{\lambda,N}\left(\frac \pi 2\right)\right)^{-1} \geq \cos(\frac \pi 3 + \frac 1{32 \pi} (\lambda-\lambda_0))^{-1} \]
	and
	\[r(K_{\lambda,N}) \leq \cos(0 - h_{\lambda,N}(0)) \leq \cos(\frac \pi 3 - \frac 5{96 \pi} (\lambda-\lambda_0))^{-1}.\]
	Estimating the difference $R(K_{\lambda,N}) - r(K_{\lambda,N})$ by its first Taylor expansion, we obtain the result for sufficiently small  $r_N>0$.
	A similar estimate holds for $\lambda \in (\lambda_0 - r_N,\lambda_0]$.
\end{proof}

\end{document}
